The development and execution of the simulation environment for the drone and car multi-agent system relied on several key software tools and libraries. Below is the list of the primary tools utilized:

\begin{itemize}
    \item \textbf{Technologies \& Software:} Visual Studio Community 2022, Git, Github, LaTeX
    \item \textbf{Platforms:} Desktop Platforms
    \item \textbf{Programming Languages:} C++, Python
    \item \textbf{Frameworks:} Unreal Engine 5.5.4, Cosys-AirSim v3.3
    \item \textbf{Database:} None required
    \item \textbf{Cloud:} None required
\end{itemize}
\noindent
The following section details the specific roles and implementation of each tool within the project workflow:

\begin{enumerate}
    \item \textbf{Cosys-AirSim (v3.3.0)}: An advanced, open-source simulator for autonomous vehicles, developed as an extension of Microsoft AirSim by the Cosys-Lab at the University of Antwerp. It is specifically designed for research-focused simulations, offering enhanced support for sensors and multi-agent systems within Unreal Engine 5.5.
    \item \textbf{Unreal Engine 5.5.4}: A high-performance 3D creation tool used to build the simulation environment. It provides the physics engine (Chaos Physics), photorealistic rendering, and the overall framework for vehicle interaction and environmental dynamics.
    \item \textbf{Python}: The primary programming language used to develop the control scripts.
    \item \textbf{C++}: The core programming language used by Unreal Engine and for the development of the Cosys-AirSim plugin.
    \item \textbf{cosysairsim package}: The official Python client wrapper for the simulation. It provides a Python interface to AirLib, facilitating communication between external control scripts and the Unreal Engine environment via RPC (Remote Procedure Call).
    \item \textbf{Visual Studio Community 2022}: The Integrated Development Environment (IDE) used for building the C++ components of the plugin and managing the Unreal Engine project dependencies.
    \item \textbf{\LaTeX}: The document preparation system used for the professional typesetting of this project report.
\end{enumerate}